% 1. 基本設定
\documentclass[unicode, aspectratio=169, 12pt]{beamer}



% 2. 日本語設定
% fontspecは指定せず、標準の原ノ味フォント（TeX Live標準）に任せるのが最も安全です
\usepackage{luatexja} 



% 3. 数式・物理系のパッケージ
\usepackage{amsmath, amssymb, mathtools, bm}
\usepackage{physics} 
% physicsとの競合（\qty）を避ける設定で読み込みます
\usepackage[overwrite-functions=false]{siunitx} 
\usepackage{tensor} % テンソルをいい感じに描くためのパッケージ，それぞれのpcで使えるか要確認！！！！！！！！



% 4. 図・グラフ・枠組み（TikZ & tcolorbox）
\usepackage{graphicx}
\usepackage{tikz}
\usetikzlibrary{math, arrows.meta, angles, quotes, decorations.markings}
\usepackage[most]{tcolorbox} % 枠付きボックス（ascmacの代わりになります）



% 5. レイアウト設定
\renewcommand{\baselinestretch}{1.2} % スライドでは1.5は広すぎるため1.2程度が一般的です
\usepackage{url}

\newcommand{\diff}[2]{\frac{\mathrm{d} #1}{\mathrm{d} #2}} %常微分を書くパッケージ
\newcommand{\pdiff}[2]{\frac{\mathrm{\partial} #1}{\mathrm{\partial} #2}} %常微分を書くパッケージ

% 6. Beamerのデザインテーマ（Boadillaベース）
\usetheme{Boadilla}
\usecolortheme{whale}
\usefonttheme{professionalfonts}
\setbeamertemplate{navigation symbols}{} % 右下のアイコンを消す
\setbeamerfont{frametitle}{size=\Large, series=\bfseries}
\setbeamersize{text margin left=1cm, text margin right=1cm} %スライドの両端1cmを白紙にして寄りすぎないようにする

% 7. タイトル情報
\title{オナーセミナー最終報告}
\author{笹伶夷\and 下橋以紗也\and 田處裕哉}
\institute{大阪大学 理学部 物理学科}
\date{\today}


%===============================================================
% 表記のルールの例
% ベクトル空間は $\bm{V}$
% ベクトルは太字にせずにそのまま x \in \bm{V} のように描く
% テンソルはtensorパッケージを使って  \tensor{t}{^\mu_nu}のように書く (このパッケージじゃないと添字の位置がうまく表示されない)
% 常微分は \diff{f}{t}のように書く, n階微分は \diff{^n f}{t^n}のように書く
% 偏微分は \pdiff{f}{t}のように書く, n階微分は \pdiff{^n f}{t^n}のように書く


%===============================================================
\begin{document}

%タイトルスライド
\begin{frame}
    \titlepage
\end{frame}


%目次スライド
\begin{frame}
    \tableofcontents
\end{frame}



%===============================================================
\section{場の古典論とは}


\begin{frame}{学部で習う物理の分野} %ここは突っ込まれそうなので直したい
    \begin{itemize}
        \item 力学 : 物体の運動を扱う
        \item 電磁気学 : 電磁場を扱う
        \item 特殊相対性理論 : 平坦な時空の物理を扱う
        \item 一般相対性理論 : 重力で歪んだ時空の物理を扱う
        \item 熱力学 : マクロな系を扱う
        \item 統計力学 : マクロな系を統計的に扱う
        \item 量子力学 : ミクロな系を扱う
    \end{itemize}
\end{frame}


\begin{frame}{場の古典論とは}

    \begin{block}{}
        場 : 空間に分布している物理量のこと\\
        古典論 : 量子力学的ではないという意味
    \end{block}

    \begin{center}
        {\Large $\Downarrow$}
    \end{center}

    \begin{alertblock}{場の古典論}
        量子力学的ではない，\\
        空間に分布する物理量(質点や電磁場，重力)を統一的に扱う分野
    \end{alertblock}

    本発表では質点と電磁場の理論について紹介する

\end{frame}



%===============================================================
\section{特殊相対性理論の舞台}

\begin{frame}{特殊相対性理論の舞台}
    \begin{block}{特殊相対性理論の出発点}
        \begin{itemize}
            \item 宇宙の情報の伝達速度には上限がある
            \item 
        \end{itemize}
    \end{block}
\end{frame}

\begin{frame}{特殊相対性理論の舞台}
    \begin{block}{ニュートン力学の舞台}
        空間３次元の実ベクトル空間$\bm{V}$を考える
        \begin{itemize}
            \item  空間上の位置を位置ベクトル$x\in\bm{V}$で表す：
            \begin{align*}
                x = (x^i) = (x,y,z)^T\quad(i=1,2,3)
            \end{align*}
            \item 時刻$t$は単なる実数パラメータで質点の運動は$\bm{r}(t)\in\bm{V}$と表す。
            \item ベクトル$a\in\bm{V}$の"大きさ"を計量$I=\mathrm{diag}(1,1,1)$を用いて定義する：
            \begin{align*}
                \norm{a}^2=a^i a^j \tensor{I}{_i_j} = (a_1)^2+(a_2)^2+(a_3)^2
            \end{align*}
        \end{itemize}        
    \end{block}
\end{frame}


\begin{frame}{特殊相対性理論の舞台}
    \begin{block}{ミンコフスキー空間 $\bm{M}$}
        時間 1 次元と空間 3 次元を統合した 4 次元実ベクトル空間 $\bm{M}$ を考える。
        \begin{itemize}
            \item 位置と時間をまとめて \textbf{4元位置ベクトル} $x \in \bm{M}$ で表す：
            \begin{equation*}
                x = (x^\mu) = (ct, x, y, z)^T \quad (\mu = 0, 1, 2, 3)
            \end{equation*}
            \item 質点の運動はミンコフスキー空間上の曲線として表される。
        \end{itemize}
    \end{block}
\end{frame}

\begin{frame}{特殊相対性理論の舞台}
    \begin{block}{ミンコフスキー空間$\bm{M}$}
        \begin{itemize}
            \item ベクトル $a, b \in \bm{M}$ の"大きさ"を計量 $g = \mathrm{diag}(1, -1, -1, -1)$ を用いて定義する：
            \begin{equation*}
                \norm{a}^2 = a^\mu a^\nu \tensor{g}{_\mu_\nu} = (a_0)^2-(a_1)^2-(a_2)^2-(a_3)^2
            \end{equation*}
        \end{itemize}        
    \end{block}
\end{frame}

\begin{frame}{特殊相対性理論の舞台}
    \begin{block}{ミンコフスキー空間$\bm{M}$}
        \begin{itemize}
            \item $g$ により、$\bm{M}$の元である共変ベクトルとその双対空間 $\bm{M}^*$ の元である双対ベクトルに1対1の対応関係を作ることができる：
            \begin{equation*}
                  a_\mu = g_{\mu\nu} a^\nu = (a_0, a_1, a_2, a_3)
            \end{equation*}
            \item 物理的な対象は$\bm{M}$と$\bm{M}^*$やそれらのテンソル積を用いて表すことができる，物理では係数のみを書いて計算していくことが多い。
        \end{itemize}
    \end{block}
\end{frame}

\begin{frame}{特殊相対性理論の舞台}
    \begin{block}{ミンコフスキー空間$\bm{M}$}
        \begin{itemize}
            \item 慣性系同士の4元位置ベクトル$x^\mu$はローレンツ変換テンソル$\Lambda=\tensor{\Lambda}{_\mu_\nu}$を用いて
            \begin{align*}
                x'^\mu = x^\nu \tensor{\Lambda}{^\mu_\nu}
            \end{align*}
            として計算できる。他の反変ベクトルや共変ベクトル，テンソルも同様に変換される。
            \item ローレンツ変換テンソルは計量$g$を不変に保つ：
            \begin{align*}
                \tensor{g}{_\mu_\nu}=\tensor{g}{_\rho_\sigma}\tensor{\Lambda}{^\rho_\mu}\tensor{\Lambda}{^\sigma_\nu}
            \end{align*}
        \end{itemize}
    \end{block}    
\end{frame}

%===============================================================
\section{最小作用の原理とは}
\begin{frame}{最小作用の原理とは}
    \begin{block}{ラグラジアン$L$}
        系の状態によって定まるエネルギーの次元を持つ量である\\
    \end{block}
ニュートン力学では$L(x,v)=\frac{1}{2}m v^2 -U(x)$と書けることが多い。
    \begin{block}{作用$I$}
        系の各瞬間のラグラジアンの値を時間積分した値：
        \begin{align*}
            I = \int_{t_0}^{t_1} L \>\mathrm{d}t
        \end{align*}
        系の時間発展が定まると作用$I$の値も一意に定まる
    \end{block}
\end{frame}


\begin{frame}{最小作用の原理とは}
    \begin{block}{最小作用の原理}
        系が時間発展するとき，実現する時間発展の作用の変分
        \footnote{作用の変分:考える系の時間発展と，少しだけ変化させた系の時間発展での作用の差のこと}
        は$\delta I=0$となる。
    \end{block}
ニュートン力学の場合は最小作用の原理から運動方程式(オイラーラグランジュ方程式)を導出できる。
\begin{align*}
    0 = \delta I
    &= \int_{t_0}^{t_1}L(x+\delta x, v + \delta v)\>\mathrm{d}t - \int_{t_0}^{t_1} L(x,v) \>\mathrm{d}t \\
    &= \int_{t_0}^{t_1} \delta x \left[\diff{}{t}\pdiff{L}{v}-\pdiff{L}{x}\right]\>\mathrm{d}t
\end{align*}
\end{frame}

\begin{frame}{最小作用の原理とは}
    \begin{block}{作用，ラグラジアンの選び方}
        最小作用の原理の結果，現実世界を再現できるようなものであるなら作用はどのようなものを選んでも良い。\\
        一般な系決める方法は存在しないが，
        \begin{itemize}
            \item 系の対称性
            \item 作用の次元が$\mathrm{ML^2T^{-1}}$になる
            \item 作用がローレンツスカラーである\footnote{相対性原理から作用がローレンツスカラーであることが導ける}
        \end{itemize}
        などの要求からある程度の任意性を除いてラグラジアンや作用を定めることができる。 %ここは怪しい
    \end{block}
\end{frame}



%===============================================================
\section{質点の時間発展}
\begin{frame}{質点の時間発展}
    自由質点の時間発展を表す作用がどのようになるのかを考える。
\end{frame}


%===============================================================
\section{電磁場の時間発展}
\begin{frame}{質点の時間発展}
    
\end{frame}

\end{document}
